\documentclass[11pt]{article}
\usepackage[margin=1in]{geometry}
\usepackage{amsmath,amssymb,amsthm,mathtools}
\usepackage{graphicx}
\usepackage{booktabs}
\usepackage{enumitem}
\usepackage{tikz}
\usetikzlibrary{positioning,arrows.meta}
\usepackage[colorlinks=true,linkcolor=blue,citecolor=blue,urlcolor=blue]{hyperref}

\newtheorem{theorem}{Theorem}
\newtheorem{corollary}[theorem]{Corollary}
\newtheorem{lemma}[theorem]{Lemma}
\newtheorem{proposition}[theorem]{Proposition}
\theoremstyle{definition}
\newtheorem{definition}[theorem]{Definition}
\newtheorem{remark}[theorem]{Remark}
\newtheorem{conjecture}[theorem]{Conjecture}

\newcommand{\N}{\mathbb{N}}
\newcommand{\cP}{\mathcal{P}}
\newcommand{\NOR}{\operatorname{NOR}}
\newcommand{\low}{\operatorname{low}}
\newcommand{\Run}{\operatorname{Run}}
\newcommand{\nmoves}{38}
\newcommand{\nclocks}{30}

\title{Subtraction Games Can Count and Compute:\\ Exponentially Long Periods, a Counterexample to Guy's
Conjecture, and Universality with a Bounded Number of Moves}
\author{}
\date{September 2026}

\begin{document}
\maketitle

\begin{abstract}
In the subtraction game with a finite move set $S\subset\N_{>0}$, a heap of $n$ tokens is reduced
by some $s\in S$ per move and the player who cannot move loses.  The sequence of winning and losing
positions is eventually periodic, and a long-standing question, going back to R.~K.~Guy, asks how
long the period can be.  Guy conjectured that the period is bounded by a polynomial of degree
$\binom{n}{2}$ in the largest move when $|S|=n$; Flammenkamp (1997) conjectured, from experiments,
that exponential periods occur.  Whether super-polynomial periods are possible at all has remained
open (Miklós--Post 2024, Larsson--Saha 2024, Larsson--Manabe 2026).

We settle the question.  For every $k\ge2$ we give an explicit set $S_k$ of exactly $\nmoves$ moves
with $\max S_k=377k-13$ whose outcome sequence has least period exactly
\[377\cdot k\cdot 2^k=(\max S_k+13)\cdot 2^{(\max S_k+13)/377}.\]
So periods of subtraction games with a bounded number of moves grow exponentially in the largest
move: Guy's conjecture is false, for outcome sequences and for nim-sequences, and Flammenkamp's
conjecture is true.  The game is a nonlinear feedback shift register whose feedback is a NOR gate;
by laying out clocks, wires and covering residues in blocks of length $377$ we make it run a
bit-serial binary counter.  The layout was found by computer search (Sidon-type constructions,
simulated annealing and exact constraint programming, reducing the number of moves
$1061\to169\to54\to\nmoves$), and every step of the proof is formally verified in Lean~4 with
Mathlib.

We then show that one-heap subtraction games are computationally universal \emph{even with a
bounded number of moves}: for every cellular automaton $F$ there is a fixed list of move templates
whose instances $S_L$, one for each ring length $L$, simulate $F$ on every initial configuration.
Consequently, deciding whether a heap size $n$ (in binary) is a P-position is PSPACE-complete already
for games with a bounded number of moves written in unary---indeed for a single explicit
one-parameter family of games---and EXPTIME-complete for a bounded number of moves written in binary,
so that in the binary setting no polynomial-time algorithm exists, unconditionally.  The layout lemma
behind these results is also formally verified.
\end{abstract}

\section{Introduction}

Let $S\subset\N_{>0}$ be finite.  In the \emph{subtraction game} $S$ a position is a heap of
$n\ge 0$ tokens; a move removes $s\in S$ tokens provided $s\le n$; under normal play the player
unable to move loses.  A position is a \emph{P-position} (a win for the Previous player) iff every
move leads to a non-P-position:
\begin{equation}\label{eq:P}
  n\in\cP_S \iff \bigl(n-s\notin \cP_S \text{ for every } s\in S \text{ with } s\le n\bigr).
\end{equation}
Equivalently, $\cP_S$ is the \emph{greedy} subset of $\N$ containing no two elements whose
difference lies in $S$.  Writing $w(n)=1$ if $n\in\cP_S$ and $w(n)=0$ otherwise, \eqref{eq:P} reads
\[ w(n)=\NOR\bigl(w(n-s): s\in S\bigr), \qquad w(\text{negative})=0. \]
Thus $w$ is generated by a shift register of length $m=\max S$ with a NOR feedback function.  By the
pigeonhole principle it is eventually periodic, with period at most $2^m$.

Subtraction games are the first problem in Nowakowski's \emph{Unsolved Problems in Combinatorial
Games}~\cite{Nowakowski}:
\begin{quote}
``In general, period lengths can be surprisingly long, and it has been suggested that they could
be super-polynomial in terms of the size of the subtraction set.  However, Guy conjectures that they
are bounded by polynomials of degree at most $\binom{n}{2}$ in $s_n$, the largest member of a
subtraction set of cardinality $n$.''
\end{quote}
Althöfer and Bültermann~\cite{AB} found families with cubic period and asked whether
super-polynomial periods occur for $|S|\ge 5$.  In his thesis, Flammenkamp~\cite{Flammenkamp}
computed record periods of ``max-symmetric'' five-move sets and observed growth like
$2^{0.3\max S}$; his conjecture that exponential periods exist is Conjecture~1 in the survey of
Larsson and Saha~\cite{LS} and is restated as open by Larsson and Manabe in 2026~\cite{LM}.
Miklós and Post~\cite{MP} write that ``it is a long time open question if exponential period length
is possible''; they proved super-polynomial periods $e^{\Omega(\sqrt{\max S})}$ for a
\emph{modified} game in which the outcomes of the first positions are prescribed by a seed.  For the
actual game no super-polynomial example was known.

\subsection{Main results}

\begin{theorem}[Exponential periods]\label{thm:main}
There is an explicit family of sets $S_k\subset\N_{>0}$ $(k\ge 2)$ such that
\begin{enumerate}[label=(\roman*)]
\item $|S_k|=\nmoves$ for every $k$: the $35$ moves $5,7,8,\dots,547$ listed in Section~\ref{sec:design}
together with $377(k-1)+222$, $377(k-1)+259$ and $377(k-1)+364$;
\item $\max S_k=377k-13$;
\item the outcome sequence of the subtraction game $S_k$ has least (eventual) period exactly
\[P_k=377\cdot k\cdot 2^{k}=(\max S_k+13)\cdot 2^{(\max S_k+13)/377};\]
in particular $P_k> 2^{\max S_k/377}$.  The nim-sequence of $S_k$ has a least period divisible by
$P_k$.
\end{enumerate}
\end{theorem}

\begin{corollary}[Guy's conjecture fails]\label{cor:guy}
For every $C,d\ge 1$ there is a subtraction set $S$ with exactly $\nmoves$ moves whose outcome period
(and nim-period) exceeds $C\,(\max S)^d$.  Concretely, $S_{15800}$ has $\nmoves$ moves, largest move
$5\,956\,587$, and period $377\cdot 15800\cdot 2^{15800}>(\max S)^{\binom{38}{2}}$, contradicting
Guy's bound for $n=\nmoves$; $k=15800$ is the least $k$ for which this happens.
\end{corollary}

Adding moves that are ``invisible'' to the construction (Remark~\ref{rem:pad}) gives the same
conclusion for every $n\ge\nmoves$ moves.

\begin{theorem}[Universality with a bounded number of moves; Section~\ref{sec:univ}]\label{thm:univintro}
For every cellular automaton $F$ with finitely many states there are a block length $B$ and a
finite list of move templates $(a_t,c_t,\delta_t)$ such that, for every ring length $L\ge3$, the
subtraction game with moves $S_L=\{(a_tL+c_t)B+\delta_t\}$ runs $F$ on a ring of $L$ cells, one after
another, on \emph{every} initial configuration: the outcome of position $((j2^L+T)L+i)B+\rho$ is a
bit of the state of cell $i$ after $T$ steps from the $j$-th configuration.
\end{theorem}

\begin{theorem}[Complexity; Section~\ref{sec:complexity}]\label{thm:complexintro}
Deciding whether $n$ (in binary) is a P-position of the subtraction game $S$ is
\begin{enumerate}[label=(\alph*)]
\item PSPACE-complete when $S$ is given in unary, even if $|S|\le c$ for a fixed constant $c$; in
fact there is one explicit family $(S_L)_{L\ge3}$ with $|S_L|\le c$ for which the problem
``given $L$ in unary and $n$ in binary, is $n\in\cP_{S_L}$?'' is PSPACE-complete;
\item EXPTIME-complete when $S$ is given in binary, even if $|S|\le c$.  Since
$\mathrm{P}\ne\mathrm{EXPTIME}$, no polynomial-time algorithm decides the outcome of subtraction
games with $c$ moves written in binary.
\end{enumerate}
\end{theorem}

Theorem~\ref{thm:complexintro} contrasts with the fact that for every \emph{fixed} game the question
is trivial (compute the period once), and with games of two moves, whose periods have closed forms.

\subsection{What exactly is proved}\label{sec:scope}

Because the question has several natural variants, we state precisely which one Theorem~\ref{thm:main}
answers.
\begin{itemize}[leftmargin=*]
\item \textbf{The standard game.}  One heap, a finite set $S$ of positive integers, normal play, and
outcomes defined for every $n\ge0$ by the standard recursion~\eqref{eq:P}.  No seed or prescribed
initial outcomes (as in~\cite{MP}), no modified starting configuration, no additional heaps, no
infinite move set, no change of the winning condition.
\item \textbf{A bounded number of moves.}  $|S_k|=\nmoves$ for \emph{every} $k$; only three moves depend on
$k$.  Guy's conjecture concerns a fixed cardinality $n$, and the theorem gives it for $n=\nmoves$.
\item \textbf{The standard notion of period.}  $P_k$ is the least $p>0$ with $w(n+p)=w(n)$ for all
sufficiently large $n$ (the preperiod is at most $377$), for the outcome sequence; for the
nim-sequence the least period is a multiple of $P_k$.
\item \textbf{Exponential in the largest move.}  The growth is not merely exponential in an auxiliary
parameter: $P_k=(m+13)2^{(m+13)/377}$ where $m=\max S_k$, so $\log_2 P_k\ge m/377$.
\item \textbf{Exact.}  The least period is determined exactly, not only bounded from below.
\item \textbf{Verified.}  All of Theorem~\ref{thm:main} and Corollary~\ref{cor:guy} is proved in Lean
from the definition of the game; the game was also simulated directly for $k\le 8$.
\item \textbf{Not claimed.}  We do not determine the least $n$ for which $n$-move games have
super-polynomial periods (open for $3\le n\le \nmoves-1$; Flammenkamp's data suggest $n=5$), nor the best
constant in the exponent.
\end{itemize}

\paragraph{Related work.}  Larsson and Wästlund~\cite{LW} and Larsson~\cite{L12} showed that
subtraction games on \emph{two or more} heaps emulate one-dimensional cellular automata (a ``time
heap'' and a ``tape heap''), yielding undecidability results; Gurvich and Vyalyi~\cite{GV} proved
PSPACE-hardness and exponential lower bounds for fixed multidimensional games, and related games
played on graphs are PSPACE-complete~\cite{DGPS}.  For one heap, where the outcome sequence is always
eventually periodic, we are not aware of any hardness result; recent work on one-heap games gives
closed forms for special families~\cite{LM,LS}.  Our constructions pack space and time into a
\emph{single} heap by cutting it into blocks.

\paragraph{Formal verification.}  Theorem~\ref{thm:main}, Corollary~\ref{cor:guy}, the nim-sequence
statement and the General Layout Lemma are proved in Lean~4 with Mathlib (Section~\ref{sec:lean}), with
main statements
\begin{itemize}
\item \texttt{SubGame.Concrete.exponential\_period} and \texttt{period\_eq\_maxMove} (Theorem~\ref{thm:main}),
\item \texttt{SubGame.Concrete.guy\_fails} and \texttt{not\_polynomially\_bounded} (Corollary~\ref{cor:guy}),
\item \texttt{SubGame.Concrete.nim\_period\_dvd} (nim-sequences),
\item \texttt{SubGame.GLL.simulate} (General Layout Lemma);
\end{itemize}
{\sloppy The command \texttt{\#print axioms} reports only \texttt{propext}, \texttt{Classical.choice} and \texttt{Quot.sound}.
The complexity results are proved on paper on top of the formal layout lemma; the universal circuits
were checked by exhaustive logic simulation and by simulating the actual games.\par}

\subsection{Idea of the construction}

The recursion $w(n)=\NOR(w(n-s):s\in S)$ applies the \emph{same} NOR gate at every position.  We
cut $\N$ into blocks of length $B$ and give every residue $r\in\{0,\dots,B-1\}$ a role:
\begin{itemize}
\item \emph{clocks} $r\in C$: $w(\tau B+r)=1$ for all blocks $\tau$;
\item \emph{zero residues}: $w(\tau B+r)=0$, forced by a move landing on a clock;
\item \emph{wires} $r=\rho(j)$: the value $w(\tau B+\rho(j))$ is a logical signal $\ell_j(\tau)$.
\end{itemize}
A move $s=qB+\delta$ applied at a wire at residue $e$ reads residue $(e-\delta)\bmod B$ in an earlier
block.  If the residues are arranged so that each wire ``sees'' exactly its intended inputs and never a
clock, and clocks see nothing but zero residues, then every wire computes a NOR gate of chosen
inputs with chosen delays (Theorem~\ref{thm:sim}).  NOR gates are universal, and a move of length
about $kB$ gives a delay line of $k$ blocks.  We use this to build a bit-serial binary counter
whose $k$-bit word is stored in the delay line: the game counts from $0$ to $2^k-1$ over and over,
and the period is exactly $k\cdot 2^k$ blocks.  Replacing the counter by a cellular automaton whose
cells are processed one per block gives the universality results.

